\documentclass[11pt,letterpaper]{report}
\newcommand{\Empresa}{Synaptix}
\newcommand{\EmpresaArchivo}{SYNAPTIX}
\newcommand{\NombreOferta}{Puerto Deportivo Bahía Panitao}
\newcommand{\VersionDocumento}{0.1}
\newcommand{\FechaDocumento}{[PENDIENTE]}
\newcommand{\ClasificacionDocumento}{Confidencial}

\usepackage{estilo/propuesta-tecnica}
% Portada vectorial reutilizable. No requiere una imagen de fondo.
\NewDocumentCommand{\PortadaPropuesta}{m m}{%
  \begin{titlepage}
    \thispagestyle{empty}
    \begin{tikzpicture}[remember picture,overlay]
      \fill[SynNavy] (current page.south west) rectangle (current page.north east);
      \fill[SynNavyLight,opacity=.52]
        ($(current page.north east)+(-4.2cm,0)$) rectangle (current page.south east);

      % Red de conexiones inspirada en la identidad digital de Synaptix.
      \draw[SynCyan,line width=2.2pt,opacity=.92]
        ($(current page.south west)+(8.2cm,2.2cm)$)
        .. controls +(3.4cm,4.0cm) and +(-4.4cm,-3.0cm) ..
        ($(current page.north east)+(-1.0cm,-4.2cm)$);
      \draw[SynBlue,line width=1.2pt,opacity=.78]
        ($(current page.south west)+(10.0cm,0.4cm)$)
        .. controls +(1.0cm,5.7cm) and +(-3.3cm,-1.0cm) ..
        ($(current page.east)+(0cm,3.0cm)$);
      \draw[SynTeal,line width=1.25pt,opacity=.70]
        ($(current page.south west)+(6.4cm,5.0cm)$)
        .. controls +(4.2cm,2.4cm) and +(-2.2cm,-4.0cm) ..
        ($(current page.north east)+(-2.0cm,-1.0cm)$);
      \draw[SynCyan!65,line width=.8pt,opacity=.62]
        ($(current page.south west)+(11.4cm,3.0cm)$)
        .. controls +(2.6cm,2.8cm) and +(-1.5cm,-3.2cm) ..
        ($(current page.north east)+(-.4cm,-7.8cm)$);

      \foreach \p/\s in {
        {($(current page.south west)+(8.2cm,2.2cm)$)}/3.2pt,
        {($(current page.south west)+(11.0cm,6.2cm)$)}/4.2pt,
        {($(current page.east)+(-1.3cm,3.0cm)$)}/3.2pt,
        {($(current page.north east)+(-2.0cm,-1.0cm)$)}/4.4pt,
        {($(current page.north east)+(-1.0cm,-4.2cm)$)}/5.0pt,
        {($(current page.north east)+(-.4cm,-7.8cm)$)}/3.1pt
      }{
        \fill[SynCyan] \p circle (\s);
        \draw[SynCyan!45,line width=.7pt] \p circle ({\s+3.2pt});
      }

      \node[anchor=north west] at ($(current page.north west)+(1.55cm,-1.35cm)$)
        {\fontsize{18}{20}\selectfont\textbf{\textcolor{SynWhite}{SYNAP}\textcolor{SynCyan}{TIX}}};
      \node[anchor=north east,text=SynCyan] at ($(current page.north east)+(-1.45cm,-1.46cm)$)
        {\fontsize{9}{11}\selectfont\bfseries INTEGRACIÓN · DATOS · DECISIONES};

      \node[anchor=west,align=left,text width=.66\paperwidth]
        at ($(current page.west)+(1.55cm,2.7cm)$) {%
          {\fontsize{13}{16}\selectfont\bfseries\textcolor{SynCyan}{PROPUESTA}}\\[2mm]
          {\fontsize{34}{38}\selectfont\bfseries\textcolor{SynWhite}{TÉCNICA}}\\[7mm]
          {\fontsize{12}{15}\selectfont\textcolor{SynWhite!82}{\NombreOferta}}
        };

      \node[anchor=west,align=left,text width=.68\paperwidth]
        at ($(current page.west)+(1.55cm,-1.0cm)$) {%
          {\fontsize{10}{12}\selectfont\bfseries\textcolor{SynCyan}{SUBDOCUMENTO #1}}\\[2mm]
          {\fontsize{19}{23}\selectfont\bfseries\textcolor{SynWhite}{#2}}
        };

      \draw[SynCyan,line width=1.5pt]
        ($(current page.south west)+(1.55cm,3.25cm)$) -- ++(3.2cm,0);
      \node[anchor=south west,align=left,text=SynWhite!74]
        at ($(current page.south west)+(1.55cm,1.35cm)$) {%
          \fontsize{9}{12}\selectfont
          \textbf{Versión} \VersionDocumento\quad
          \textbf{Fecha} \FechaDocumento\quad
          \textbf{Clasificación} \ClasificacionDocumento\\[2mm]
          \textbf{Firmante autorizado:} \FirmanteDocumento
        };
    \end{tikzpicture}
    \null
  \end{titlepage}
  \clearpage
}


\newcommand{\CampoFicha}[2]{#1 & #2 \\ \hline}
\newcommand{\FichaInnovacion}[3]{%
  \section*{Ficha #1: #2}
  \addcontentsline{toc}{section}{Ficha #1: #2}
  La Tabla~\ref{tab:t19-#1} presenta los diecisiete campos del Formulario T-19 para la innovación #1.
  \begin{longtable}{L{\dimexpr.30\textwidth-2\tabcolsep\relax}L{\dimexpr.70\textwidth-2\tabcolsep\relax}}
    \caption{Ficha #1 de innovación: #2}\label{tab:t19-#1} \\
    \hline
    \EncabezadoTabla{Campo} & \EncabezadoTabla{Contenido} \\
    \hline
    \endfirsthead
    \hline
    \EncabezadoTabla{Campo} & \EncabezadoTabla{Contenido} \\
    \hline
    \endhead
    #3
  \end{longtable}
  \FuenteTabla{Antecedentes y diseño de Synaptix para la innovación #1; fuentes identificadas en la ficha.}
}

\begin{document}
\pagenumbering{arabic}
\renewcommand{\thetable}{T-19.\arabic{table}}
\PortadaFormulario{T-19}{Cartera de innovaciones}
\tableofcontents
\clearpage
\chapter*{Fichas de innovación}
\addcontentsline{toc}{chapter}{Fichas de innovación}

Synaptix presenta una ficha por cada tipo del artículo 28. Los campos siguen el Formulario T-19 y los siete elementos del artículo 29 \parencite[pp.~19--20 y 64]{bases-admin}. Las fichas se corresponden con los apartados 13.1 a 13.5 del SD13. Los valores de línea base que no aporta el caso requieren medición; las metas son propuestas de diseño. Los códigos de las extensiones 4 y 5 se conciliarán con los catálogos de SD4 y SD7. La valorización monetaria se presenta en la Oferta Económica, conforme al artículo 50.2 \parencite[p.~30]{bases-admin}.

%Reemplaza la ficha 1 por este bloque. Incluye el calendario, la EDT, los indicadores y los riesgos de la última versión, sin estimaciones de horas-hombre y con las contingencias explícitas. Subdocumento_13___SYNAPTIX-5.pdf

\FichaInnovacion{1}{Servicio de Solicitud y Seguimiento de Mantenimiento de Embarcaciones}{%

\CampoFicha{Tipo de innovación}{%
Producto o servicio de base tecnológica.
}

\CampoFicha{Nombre}{%
Servicio de Solicitud y Seguimiento de Mantenimiento de Embarcaciones.
}

\CampoFicha{Problema u oportunidad}{%
Dar trazabilidad a las solicitudes de mantenimiento entre armador,
marina y contratista, desde su envío hasta su cierre. Complementa
la autorización de acceso asociada a la orden de trabajo de
RF-CTR-04 con seguimiento del avance e información para el armador.
}

\CampoFicha{Tecnología o práctica}{%
Flujo digital de solicitudes con estados, responsables, historial
y notificaciones. La marina revisa y deriva; el contratista acepta
o rechaza y actualiza el avance; el armador consulta el estado.
La asignación no es automática y el cierre administrativo no
certifica la calidad técnica de la reparación.
}

\CampoFicha{Madurez y escala}{%
TRL 2 preliminar para el servicio integrado: concepto y flujo
definidos, aún sin implementación ni validación en la marina.
El avance de madurez dependerá de las pruebas y del piloto.
}

\CampoFicha{Fuentes APA 7}{%
IBM (s.f.-a, s.f.-b): creación y estados de órdenes de trabajo;
Manning (2023): escala TRL; Pontificia Universidad Católica de Valparaíso
(2026): contexto operacional. Referencias completas en SD13.
}

\CampoFicha{Inserción en arquitectura}{%
Componente IN1-SOL en Varadero y contratistas. Interfaces
IN1-IF-01 a IN1-IF-06 con AUT, CTR, VAR, ALE, ADM y Clientes
y contratos para solicitudes, habilitación, órdenes, avisos,
permisos y vínculo del armador. Incluye control de duplicados,
versiones e historial auditable. Ante indisponibilidad se
utilizará un registro controlado y posterior conciliación.
}

\CampoFicha{Paquetes de la EDT}{%
IN1.1.1 a IN1.1.3: diseño; IN1.2.1 a IN1.2.3: servicio y
portal; IN1.3.1 a IN1.3.3: integración; IN1.4.1 a IN1.4.2:
pruebas; IN1.5.1 a IN1.5.6: capacitación, línea base y piloto;
IN1.6.1: despliegue; IN1.7: operación y seguimiento.
Desglose en la tabla 13.1 del SD13 y correspondencia con SD7/T-15.
}

\CampoFicha{Mes de materialización}{%
Mes 21, después de la aceptación. Diseño: 14--15; desarrollo
e integración: 16--17; pruebas: 18; línea base y capacitación:
19; piloto y evaluación: 20. Operación durante los meses 21--56.
}

\CampoFicha{Inversión requerida}{%
Diseño, desarrollo, integración, pruebas, capacitación y puesta
en marcha, asociados a IN1.1 a IN1.6. Se valorizarán los recursos
incrementales en SD11 y la Oferta Económica, reutilizando los
componentes financiados por la solución base.
}

\CampoFicha{Costo operacional}{%
Soporte, mantenimiento, almacenamiento del historial,
notificaciones y asistencia a usuarios durante los meses
21--56, asociados a IN1.7. Se considerará el esfuerzo de
registro asistido y se evitará duplicar costos del servicio base.
}

\CampoFicha{Beneficio y cuantificación}{%
Reducir el esfuerzo administrativo de seguimiento, coordinación
y búsqueda de antecedentes. Se comparará el trabajo requerido
por solicitudes equivalentes antes y durante el piloto.
La capacidad liberada será un beneficio operacional; sólo se
reconocerá ahorro económico con gasto evitado u otro beneficio
valorizado y aceptado, descontando los costos incrementales.
El beneficio no se duplicará con la innovación 4.
}

\CampoFicha{Indicador, línea base y meta}{%
Mediana de horas entre envío y primera respuesta explícita del
contratista, sea aceptación o rechazo. Base: solicitudes
comparables del mes 19. Meta propuesta: reducción del 20\,\%
en el mes 20, sin disminuir la proporción respondida en siete
días. Indicador complementario: 100\,\% de transiciones con
estado anterior y nuevo, fecha y responsable.
}

\CampoFicha{Momento de medición}{%
Línea base en el mes 19 y piloto en el mes 20. Se incluirán
solicitudes ingresadas durante los primeros 21 días de cada
mes, observadas por siete días, con un mínimo operativo de
30 solicitudes por condición. Se informarán casos sin respuesta,
cancelados y excluidos. Con muestra insuficiente, base cero
o falta de comparabilidad, el beneficio será no demostrado.
}

\CampoFicha{Riesgo, probabilidad e impacto}{%
Valoraciones iniciales: R-IN1-01, actualización insuficiente
por contratistas, probabilidad alta e impacto alto;
R-IN1-02, baja participación de armadores, probabilidad media
e impacto medio; R-IN1-03, fallas de integración, probabilidad
media e impacto alto; R-IN1-04, acceso no autorizado,
probabilidad baja e impacto muy alto.
}

\CampoFicha{Mitigación}{%
R-IN1-01 y R-IN1-02: capacitación, orientación, recordatorios
y responsabilidades claras, a cargo del Líder de Implantación.
R-IN1-03: definición de interfaces y pruebas de integración,
a cargo del Arquitecto de Solución.
R-IN1-04: validación de permisos, vínculo del armador y
asignación del contratista, a cargo del Encargado de Seguridad.
}

\CampoFicha{Contingencia}{%
R-IN1-01: registro asistido por la marina si un cambio comunicado
lleva más de un día hábil sin registrarse, conservando fuente
y fecha.
R-IN1-02: reforzar la convocatoria con menos de 15 solicitudes
a mitad del piloto; si al cierre no se alcanzan 30, informar
beneficio no demostrado y proponer ampliar la observación.
R-IN1-03: postergar el piloto afectado y mantener la solución
base si falla una interfaz crítica al cierre del mes 18.
R-IN1-04: suspender la función ante acceso indebido confirmado,
preservar evidencia y corregir y repetir pruebas antes de
reactivarla. Los ajustes de calendario se tramitarán mediante
control de cambios.
}

}


\FichaInnovacion{2}{Programa de Mantenimiento Preventivo en Invernada}{%
  \CampoFicha{Tipo de innovación}{Proceso.}
  \CampoFicha{Nombre}{Programa de Mantenimiento Preventivo en Invernada.}
  \CampoFicha{Problema u oportunidad}{180 puestos de invernada, 152 ocupados en invierno de 2026; 14 puestos de trabajo del varadero; 62\,\% de los 1.150 movimientos anuales del travelift concentrados en dos campañas de seis semanas. La línea base de trabajos pendientes a la botadura no está cuantificada \parencite[caps.~2 y 7]{caso07}.}
  \CampoFicha{Tecnología o práctica}{Ficha de ingreso, propuesta voluntaria de tareas preventivas, calendario de invierno según secuencia de botadura y revisión de cierre. Cada trabajo aceptado se vincula a la solicitud de la innovación 1. El armador conserva la decisión de contratar y elegir al contratista.}
  \CampoFicha{Madurez y escala}{TRL 2 preliminar aplicado al proceso integrado para esta marina; no existe evidencia de un ciclo implantado en el caso \parencite{nasa-trl}.}
  \CampoFicha{Fuentes APA 7}{Caso 07: cifras, calendario y restricciones; UNE-EN 13306:2018: terminología de mantenimiento \parencite{une13306}; NASA: escala TRL \parencite{nasa-trl}.}
  \CampoFicha{Inserción en arquitectura}{VAR aporta posición y secuencia de la explanada; CTR, acreditación de empresas; AUT, ficha y aceptación; ALE, avisos; ADM, catálogo configurable. Integración con la innovación 1. Interfaces por conciliar con SD4.}
  \CampoFicha{Paquetes de la EDT}{IN2.1 diseño, meses 14--15; IN2.2 ficha y propuestas, 16--17; IN2.3 calendario e integración, 17--18; IN2.4 pruebas y aceptación técnica, 19--20; IN2.5 piloto y evaluación, de \(M_{\mathrm{inv}}\) a \(M_{\mathrm{inv}}+5\); IN2.6 transferencia, \(M_{\mathrm{inv}}+5\); IN2.7 operación incremental, de \(M_{\mathrm{inv}}\) a 56. Paquetes terminales en la tabla 13.3 del SD13; correspondencia de códigos y cargas por conciliar en SD7/T-15.}
  \CampoFicha{Mes de materialización}{\(M_{\mathrm{inv}}\) será el primer junio disponible con IN1 y las interfaces aceptadas, no anterior al mes 21. Piloto en junio--agosto, de \(M_{\mathrm{inv}}\) a \(M_{\mathrm{inv}}+2\); evaluación de botaduras y transferencia hasta \(M_{\mathrm{inv}}+5\). Ciclo completo desde abril--mayo siguiente, \(M_{\mathrm{inv}}+10\), si el plazo permite completar la evaluación. SD7 fijará el número contractual a partir de la fecha de inicio y las restricciones del varadero.}
  \CampoFicha{Inversión requerida}{Levantamiento, diseño, calendario, integración, pruebas, capacitación y piloto. Estimación monetaria por desarrollar en la Oferta Económica.}
  \CampoFicha{Costo operacional}{Actualización del catálogo, seguimiento de planes, soporte y notificaciones desde el inicio del piloto hasta el término del servicio. Incluir únicamente sus costos incrementales en el flujo de caja.}
  \CampoFicha{Beneficio y cuantificación}{Reprogramaciones de botadura efectivamente evitadas por trabajos pendientes y costos unitarios respaldados de coordinación y movimientos adicionales, descontando costos incrementales de implementación y operación. Una reducción de pendientes no constituye por sí sola ahorro monetario. Separar inversión, costos recurrentes y ahorros verificables; no duplicar beneficios de IN1 ni atribuir ingresos de charter o derechos de trabajo sin evidencia.}
  \CampoFicha{Indicador, línea base y meta}{Embarcaciones con trabajos pendientes en su fecha originalmente prevista de botadura / embarcaciones con al menos un trabajo contratado. Conservar esa fecha aunque se reprograme. Base: grupo comparable sin programa de la misma campaña y, cuando exista, historial propio. Meta preliminar: reducir a la mitad la proporción respecto del grupo sin programa; registrar diferencias de alcance, clima, retiros y exclusiones.}
  \CampoFicha{Momento de medición}{Cierre de la campaña de botadura posterior al piloto, en \(M_{\mathrm{inv}}+5\). Complementos: proporción de trabajos ejecutados en junio--agosto y trazabilidad del 100\,\% de trabajos aceptados mediante IN1 con estado y fecha. Sin referencia positiva o evidencia suficiente, no declarar demostrada la reducción relativa.}
  \CampoFicha{Riesgo, probabilidad e impacto}{Valoraciones iniciales: R-IN2-01, baja participación de armadores, probabilidad media e impacto alto; R-IN2-02, incumplimiento de fechas por contratistas, alta/alto; R-IN2-03, suspensión por meteorología, alta/alto; R-IN2-04, indisponibilidad de IN1 o sus interfaces, media/alto. Se revisarán con el piloto y se incorporarán al plan de riesgos.}
  \CampoFicha{Mitigación}{Participación voluntaria y capacitación a cargo de Implantación; confirmación de disponibilidad y fechas por el Líder Funcional y la administración del varadero; reserva de planificación y alertas marítimas por el Jefe de Proyecto y la marina; verificación de IN1 e interfaces por Arquitectura. El armador conserva la elección de trabajos y contratista.}
  \CampoFicha{Contingencia}{Sin grupo evaluable, ampliar convocatoria o ajustar alcance sin declarar beneficio; ante atraso, reprogramar o proponer otro contratista con aceptación del armador; ante meteorología insegura, suspender y reprogramar, acordando el tratamiento de pendientes si se agota la reserva. Sin IN1 o interfaces aceptadas, no iniciar el piloto; ante interrupción, usar el procedimiento vigente con registro controlado y conciliación posterior. No exigir contratación ni alterar maniobras durante izaje.}
}

La ficha corresponde al programa de invernada del apartado 13.2 y conserva sus meses relativos y valoraciones preliminares de riesgo. SD7/T-15 deberá conciliar códigos, cargas y calendario contractual; la Oferta Económica desarrollará la valorización.

\FichaInnovacion{3}{Lectura inteligente de documentos para el ingreso asistido de antecedentes}{%

\CampoFicha{Tipo de innovación}{%
Tecnológica: OCR y extracción asistida de datos documentales.
}

\CampoFicha{Nombre}{%
Lectura inteligente de documentos para el ingreso asistido de antecedentes.
}

\CampoFicha{Problema u oportunidad}{%
Reducir la transcripción de pólizas de embarcaciones mediante campos precompletados y revisión humana, complementando RF-ADM-01 y RF-NAV-11.
}

\CampoFicha{Tecnología o práctica}{%
Amazon Textract y reglas para extraer póliza, asegurado, aseguradora, nave y vigencias. Revisión obligatoria antes de aprobar; no acredita autenticidad ni cobertura.
}

\CampoFicha{Madurez y escala}{%
TRL 2 preliminar para la integración local: concepto definido, pendiente de implementación y validación con pólizas reales.
}

\CampoFicha{Fuentes APA 7}{%
AWS: extracción documental \parencite{in3-aws-analisis,in3-aws-practicas,in3-aws-limites}; ISO: riesgos de IA \parencite{in3-iso23894}; NASA: escala TRL \parencite{in3-nasa-trl}.
}

\CampoFicha{Inserción en arquitectura}{%
AUT recibe documentos; ADM presenta la revisión al funcionario autorizado y administra permisos bajo las políticas de Auditoría y seguridad. Clientes y contratos conserva los maestros de armador y embarcación; NAV aplica las validaciones de vigencia sobre los datos aprobados. La identidad procede de la capa de Seguridad transversal del SD4. Adaptador asíncrono conectado al repositorio y la auditoría existentes.
}

\CampoFicha{Paquetes de la EDT}{%
IN3.1 diseño, meses 14--15; IN3.2 extracción, 16; IN3.3 revisión e integración, 17; IN3.4 pruebas y matriz de controles de IA, 18; IN3.5 preparación y piloto, 18--20; IN3.6 transferencia, 21; IN3.7 operación incremental, 21--56. Paquetes terminales en la tabla 13.6 del SD13; correspondencia y cargas por conciliar con SD7/T-15.
}

\CampoFicha{Mes de materialización}{%
Mes 21, tras piloto de los meses 19--20 y aceptación. Diseño: 14--15; implementación: 16--17; pruebas y preparación: 18. Requiere interfaces de AUT, ADM, NAV, repositorio y auditoría disponibles, y controles de IA aceptados. SD7 deberá confirmar perfiles compartidos y dependencias; los despliegues respetarán las ventanas protegidas.
}

\CampoFicha{Inversión requerida}{%
Levantamiento documental, configuración del extractor, integraciones, interfaz de revisión, pruebas y capacitación, según hitos de los meses 14--21. Valorización en la Oferta Económica sin duplicar repositorio, identidad y auditoría de la solución base.
}

\CampoFicha{Costo operacional}{%
Procesamiento de páginas, almacenamiento adicional, monitoreo y mantenimiento de reglas. Consumo del piloto en meses 19--20 y operación recurrente desde el 21. Valorización según volumen medido y condiciones comerciales aplicables.
}

\CampoFicha{Beneficio y cuantificación}{%
Capacidad operacional recuperada por reducción del tiempo de ingreso y revisión frente al método manual, incluidas correcciones y casos completados manualmente. Solo se incorporará como ahorro de caja cuando reduzca costos efectivos y respaldados. El flujo de caja distinguirá inversión, costos recurrentes y ahorros verificables; ampliar a otros documentos requerirá evaluación adicional.
}

\CampoFicha{Indicador, línea base y meta}{%
Reducción de la mediana del tiempo total de ingreso y revisión, incluidas correcciones y esperas que impidan continuar la tarea. Base: grupo manual comparable; los casos derivados a ingreso manual permanecen en el grupo asistido y la latencia del extractor se registra por separado. Meta preliminar: al menos 30\,\%, sin aumentar registros aprobados con errores ni aceptar errores críticos de asociación de embarcación o vigencias; trazabilidad del 100\,\% de aprobaciones.
}

\CampoFicha{Momento de medición}{%
Ocho semanas consecutivas dentro de los meses 19--20; evaluación al cierre del mes 20. Muestra de 120 pólizas distintas: 60 con ingreso manual y 60 asistido, comparables por formato, páginas y legibilidad, alternando métodos entre funcionarios. Pueden utilizarse pólizas históricas autorizadas en pruebas sin modificar expedientes vigentes. Los documentos de configuración quedan fuera de la evaluación. QA contrastará los 120 registros con el original. Sin muestra suficiente, informar resultado inconcluso y actualizar calendario para cualquier ampliación.
}

\CampoFicha{Riesgo, probabilidad e impacto}{%
Valoraciones iniciales: R-IN3-01, extracción incorrecta por calidad o formato, probabilidad alta e impacto alto; R-IN3-02, revisión humana insuficiente, media/alto; R-IN3-03, tratamiento fuera de condiciones autorizadas, media/alto; R-IN3-04, interrupción del extractor o conectividad, media/medio; R-IN3-05, volumen o ahorro insuficiente, media/medio. Cada evento tendrá responsable, mitigación y contingencia registrados.
}

\CampoFicha{Mitigación}{%
Datos aplicará reglas, fragmentos de origen y evaluación por formato; Implantación capacitará revisores y QA comprobará aprobaciones; Seguridad verificará permisos, ubicación, retención y proveedor; Operación/SRE supervisará fallos y reintentos. El Jefe de Proyecto revisará volumen, tiempos y costos. Ampliar solo tras validar el piloto.
}

\CampoFicha{Contingencia}{%
Suspender formatos con errores críticos o recurrentes y usar ingreso manual; ante error crítico aprobado, pausar el flujo y revisar registros afectados. Bloquear envíos incompatibles con las condiciones autorizadas y aplicar respuesta a incidentes. Ante interrupción, conciliar resultados tardíos sin duplicar ni sobrescribir datos aprobados. Sin beneficio demostrado, mantener el proceso manual y ajustar o deshabilitar la extensión, sin ampliar automáticamente su alcance.
}
}

\FichaInnovacion{4}{Servicio Gestionado por Resultados}{%
  \CampoFicha{Tipo de innovación}{Modelo de negocio o de contratación.}
  \CampoFicha{Nombre}{Servicio Gestionado por Resultados.}
  \CampoFicha{Problema u oportunidad}{36 meses de operación sin TI interna y coordinación con 62 empresas contratistas. No existe un tiempo base de seguimiento de mantenimiento en el caso \parencite[caps.~2 y 7; sec.~9.11]{caso07}.}
  \CampoFicha{Tecnología o práctica}{Componente fijo para obligaciones y variable por reducción comprobada de tiempo activo de seguimiento del servicio adicional de la innovación 1. Factor proporcional a la meta, sujeto a datos, calidad, atribución y acta. E-25 admite componentes variables durante los meses 21--56 \parencite[pp.~74--75]{bases-admin}.}
  \CampoFicha{Madurez y escala}{Diseño conceptual del modelo. Hitos de validación: diseño, piloto sin efectos de pago, aceptación de evidencia y condiciones, y conciliación en operación. No se declaran los últimos tres como realizados.}
  \CampoFicha{Fuentes APA 7}{Bases Administrativas: artículos 28--30, 50.2 y E-25; Caso 07: operación; Cabinet Office: indicadores objetivos y controlables por el proveedor \parencite[sec.~5.1.5]{uk-outcomes}.}
  \CampoFicha{Inserción en arquitectura}{Se propone IN4-MET en Varadero y contratistas/Núcleo de Servicios e IN4-CON en Finanzas y cobranza/Núcleo Administrativo, sobre ECS/Fargate según SD4, DA-04 y tabla 4-10. Enlace con solicitudes y faenas de SD5; tiempos activos con captura propia. REST/OpenAPI 3.1, escrituras idempotentes y eventos versionados conforme a DA-06; PostgreSQL y evidencia S3 conforme a DA-08. AUT ofrece el canal; ADM separa preparación y aprobación. Conciliación diferida hasta sincronización y validación, sin cobro automático. SD4 recoge la ubicación de IN4-MET e IN4-CON. Contratos, entidades y retención detallados siguen pendientes en SD4/SD5 y A5.1/A5.2.}
  \CampoFicha{Paquetes de la EDT}{IN4.1 regla y medición, meses 2--4; IN4.2 instrumentación y pruebas, 14--18; IN4.3 piloto, 19--20; IN4.4 conciliación, 21--56. Incorporar identificadores y cargas a SD7/T-15.}
  \CampoFicha{Mes de materialización}{Mes 21, tras piloto de meses 19--20 y aceptación de evidencia. Requiere disponible el flujo adicional de mantenimiento.}
  \CampoFicha{Inversión requerida}{Definición de la regla, instrumentación, expediente y pruebas, representadas por \(I_4\). Valorización en la Oferta Económica.}
  \CampoFicha{Costo operacional}{\(O_{4,m}\): medición, revisión y conciliación mensual. El proveedor lo considera aun en el escenario de incentivo cero.}
  \CampoFicha{Beneficio y cuantificación}{\(H_{4,m}=N_m\max(0,\bar T_0-\bar T_m)/60\) horas estimadas recuperadas. La valorización del beneficio corresponde a la innovación que lo produce; la 4 añade costos de medición y transferencia variable, sin sumar un segundo ahorro bruto. Las horas no constituyen automáticamente ahorro de caja.}
  \CampoFicha{Indicador, línea base y meta}{Reducción de tiempo medio activo de seguimiento frente a base prospectiva con las funciones obligatorias del SD2, incluidas RF-CTR-04, RF-AUT-06 y RF-AUT-10. Registrar tarea, responsable e intervalos; excluir esperas y ejecución del mantenimiento. Comparar composición equivalente y conservar faltantes y casos abiertos. Meta propuesta: 20\,\%, sin aumentar tasa de errores ni reducir cierres en el mismo plazo. Fijar protocolo y meta antes de medir. Factor \(F_4=\min(1,\max(0,R_4/0{,}20))\); resultado no evaluable o calidad insuficiente implica cero. El reconocimiento económico requiere además beneficio neto aceptado y tope de la Oferta Económica.}
  \CampoFicha{Momento de medición}{Base al inicio del mes 19; cierre del piloto en el 20; conciliación mensual desde el 21. Mínimo propuesto: 30 solicitudes comparables por condición, con igual plazo de observación. Synaptix prepara evidencia, la Contraparte Técnica la verifica y el Administrador del Contrato tramita el estado de pago. Sin evidencia suficiente, el variable del período es cero.}
  \CampoFicha{Riesgo, probabilidad e impacto}{Atribución incorrecta o incentivos a omitir casos difíciles. Valoración preliminar: probabilidad media, impacto alto.}
  \CampoFicha{Mitigación}{Muestra y exclusiones acordadas antes de medir; revisión de cambios externos; control de errores, cierres y reaperturas; auditoría y acta del cliente. Multas contractuales independientes.}
  \CampoFicha{Contingencia}{Suspender el reconocimiento variable, corregir la medición y reevaluar. Mantener la operación y componente fijo obligatorios; incentivo cero si no se demuestra beneficio económico aceptable.}
}

La mejora contractual depende de un resultado aceptado y no de la disponibilidad obligatoria. El flujo de caja debe distinguir ahorro bruto, costos de medición y transferencia variable para evitar duplicaciones.

\FichaInnovacion{5}{Pasaporte Ambiental Náutico con Planes de Reducción}{%
  \CampoFicha{Tipo de innovación}{Experiencia de usuario, sostenibilidad o impacto social.}
  \CampoFicha{Nombre}{Pasaporte Ambiental Náutico con Planes de Reducción.}
  \CampoFicha{Problema u oportunidad}{Cero de 320 amarras con medición individual; tres brechas ambientales descritas y unas 140 personas contratistas sin capacitación acreditada. Charter con 14 embarcaciones y 22\,\% de ingresos por amarra exige certificación vigente en 2029 \parencite[caps.~1, 8, 11 y 13; sec.~7.3]{caso07}. Los controles de base no se atribuyen al aporte innovador.}
  \CampoFicha{Tecnología o práctica}{Historial propio y plan voluntario de reducción de agua: acción, responsable, plazo, evidencia y comparación posterior. Reglas deterministas, sin IA ni control automático del suministro. Prácticas de objetivos, educación, medición y detección de fugas de WaterSense \parencite{epa-watersense-work}.}
  \CampoFicha{Madurez y escala}{TRL 2 preliminar para la incorporación local, aún sin prototipo ni piloto ejecutados \parencite{nasa-trl}.}
  \CampoFicha{Fuentes APA 7}{Caso 07: ambiente y charter; EPA: WaterSense at Work; NASA: TRL; Transversales: integración, indicadores y riesgos del capítulo 26 \parencite[pp.~44--45]{bases-transversales}.}
  \CampoFicha{Inserción en arquitectura}{IN5-PLAN se propone en Ambiente y consumos/Núcleo de Servicios sobre ECS/Fargate, con PostgreSQL; IN5-HIS utiliza historia derivada Parquet/S3/Athena según SD4 y SD5. AUT presenta la vista React. Consume MediciónConsolidada y referencias de Infraestructura, Operación marítima y Clientes y contratos; ConsentimientoModificado restringe usos. REST/OpenAPI y eventos idempotentes conforme a DA-06. Reutiliza captura local y sincronización de DA-09, sin acceso directo a MQTT ni modificación de cobros. Evaluaciones pendientes durante desconexión o con evidencia incompleta. Códigos, entidades, contratos y retención aún por detallar en SD4 y A5.1/A5.2.}
  \CampoFicha{Paquetes de la EDT}{IN5.1 diseño y consentimiento, meses 14--15; IN5.2 historial, portal y pruebas, 16--18; IN5.3 piloto, 19--20; IN5.4 seguimiento, 21--56. Incorporar correspondencia y cargas a SD7/T-15.}
  \CampoFicha{Mes de materialización}{Mes 21, tras evaluación del piloto 19--20. Requiere medición y permanencias validadas, adquisición de medidores y obras del cliente conforme a SUP-04/EXC-06/EXC-10. VAC-03 y VAC-14 siguen por validar. Conciliar con SD7/T-15, inicio contractual, congelamientos RES-11, marejadas RES-12 y serie previa a temporada 2029.}
  \CampoFicha{Inversión requerida}{\(I_5\): historial, planes, vistas, consentimiento e integración y pruebas. Valorización en la Oferta Económica.}
  \CampoFicha{Costo operacional}{\(O_{5,t}\): almacenamiento, soporte, revisión de acciones y conciliación de discrepancias. Registrar sólo costo incremental.}
  \CampoFicha{Beneficio y cuantificación}{Volumen evitado estimado \(V_t=D_t(w_0-w_t)\). Escenario objetivo: \(0{,}10w_0D_t\), sin presentarlo como medición. Valorar con costo documentado del recurso y descontar operación incremental e inversión. Asignar ahorro al actor que paga el recurso, sin duplicarlo ni atribuir el ingreso completo del charter.}
  \CampoFicha{Indicador, línea base y meta}{Agua suministrada en la marina y atribuible a la embarcación, en metros cúbicos por día efectivo de permanencia. Se reutiliza la serie obligatoria de RF-CON-01 a 06 y la consulta de RF-AUT-03 del SD2; no representa consumo total durante navegación. Base propia de 28 días del mes 19; comparar 28 días del 20. Metas propuestas y fijadas antes del piloto: reducción del 10\,\% y, por separado, acción ejecutada con evidencia en el 50\,\% de inscritos. Informar datos válidos, exclusiones y resultados desfavorables; una acción no demuestra por sí sola ahorro.}
  \CampoFicha{Momento de medición}{Evaluación en mes 20 y seguimiento mensual desde el 21, con informe previo a temporada 2029. Iniciar sólo con lecturas, permanencias y accesos validados; ejecutar acciones entre las ventanas comparadas. El responsable de la marina revisa contexto, Synaptix prepara el cálculo y la Contraparte Técnica acepta evidencia. Sin base positiva, permanencia o comparación defendible, no declarar reducción.}
  \CampoFicha{Riesgo, probabilidad e impacto}{Sesgo voluntario o datos incompletos: probabilidad media, impacto alto. Baja participación: media/medio. Acceso indebido entre flotas: baja/alto. Valoraciones preliminares.}
  \CampoFicha{Mitigación}{Criterios fijados antes del piloto; comparación propia; trazabilidad y conservación de resultados desfavorables; acciones breves y voluntarias; permisos por embarcación, pruebas de aislamiento y auditoría.}
  \CampoFicha{Contingencia}{No publicar reducciones no evaluables; ampliar observación y ajustar o desactivar el plan si no demuestra utilidad. Suspender vista afectada ante incidente de privacidad y conservar medición, residuos y capacitación obligatorios.}
}

La ficha mide una reducción de agua adicional al registro y la facturación obligatorios. El porcentaje de adopción informa sobre el uso del plan; el resultado ambiental requiere consumo comparable y no acredita una certificación.

\clearpage
\printbibliography[heading=referenciassynaptix,title={Referencias}]
\end{document}
